\IfFileExists{stacks-project.cls}{%
\documentclass{stacks-project}
}{%
\documentclass{amsart}
}

% For dealing with references we use the comment environment
\usepackage{verbatim}
\newenvironment{reference}{\comment}{\endcomment}
%\newenvironment{reference}{}{}
\newenvironment{slogan}{\comment}{\endcomment}
\newenvironment{history}{\comment}{\endcomment}

% For commutative diagrams we use Xy-pic
\usepackage[all]{xy}

% We use 2cell for 2-commutative diagrams.
\xyoption{2cell}
\UseAllTwocells

% We use multicol for the list of chapters between chapters
\usepackage{multicol}

% This is generally recommended for better output
\usepackage{lmodern}
\usepackage[T1]{fontenc}

% For cross-file-references
\usepackage{xr-hyper}

% Package for hypertext links:
\usepackage{hyperref}

% For any local file, say "hello.tex" you want to link to please
% use \externaldocument[hello-]{hello}
\externaldocument[introduction-]{introduction}
\externaldocument[conventions-]{conventions}
\externaldocument[sets-]{sets}
\externaldocument[categories-]{categories}
\externaldocument[topology-]{topology}
\externaldocument[sheaves-]{sheaves}
\externaldocument[sites-]{sites}
\externaldocument[stacks-]{stacks}
\externaldocument[fields-]{fields}
\externaldocument[algebra-]{algebra}
\externaldocument[brauer-]{brauer}
\externaldocument[homology-]{homology}
\externaldocument[derived-]{derived}
\externaldocument[simplicial-]{simplicial}
\externaldocument[more-algebra-]{more-algebra}
\externaldocument[smoothing-]{smoothing}
\externaldocument[modules-]{modules}
\externaldocument[sites-modules-]{sites-modules}
\externaldocument[injectives-]{injectives}
\externaldocument[cohomology-]{cohomology}
\externaldocument[sites-cohomology-]{sites-cohomology}
\externaldocument[dga-]{dga}
\externaldocument[dpa-]{dpa}
\externaldocument[sdga-]{sdga}
\externaldocument[hypercovering-]{hypercovering}
\externaldocument[schemes-]{schemes}
\externaldocument[constructions-]{constructions}
\externaldocument[properties-]{properties}
\externaldocument[morphisms-]{morphisms}
\externaldocument[coherent-]{coherent}
\externaldocument[divisors-]{divisors}
\externaldocument[limits-]{limits}
\externaldocument[varieties-]{varieties}
\externaldocument[topologies-]{topologies}
\externaldocument[descent-]{descent}
\externaldocument[perfect-]{perfect}
\externaldocument[more-morphisms-]{more-morphisms}
\externaldocument[flat-]{flat}
\externaldocument[groupoids-]{groupoids}
\externaldocument[more-groupoids-]{more-groupoids}
\externaldocument[etale-]{etale}
\externaldocument[chow-]{chow}
\externaldocument[intersection-]{intersection}
\externaldocument[pic-]{pic}
\externaldocument[weil-]{weil}
\externaldocument[adequate-]{adequate}
\externaldocument[dualizing-]{dualizing}
\externaldocument[duality-]{duality}
\externaldocument[discriminant-]{discriminant}
\externaldocument[derham-]{derham}
\externaldocument[local-cohomology-]{local-cohomology}
\externaldocument[algebraization-]{algebraization}
\externaldocument[curves-]{curves}
\externaldocument[resolve-]{resolve}
\externaldocument[models-]{models}
\externaldocument[functors-]{functors}
\externaldocument[equiv-]{equiv}
\externaldocument[pione-]{pione}
\externaldocument[etale-cohomology-]{etale-cohomology}
\externaldocument[proetale-]{proetale}
\externaldocument[relative-cycles-]{relative-cycles}
\externaldocument[more-etale-]{more-etale}
\externaldocument[trace-]{trace}
\externaldocument[crystalline-]{crystalline}
\externaldocument[spaces-]{spaces}
\externaldocument[spaces-properties-]{spaces-properties}
\externaldocument[spaces-morphisms-]{spaces-morphisms}
\externaldocument[decent-spaces-]{decent-spaces}
\externaldocument[spaces-cohomology-]{spaces-cohomology}
\externaldocument[spaces-limits-]{spaces-limits}
\externaldocument[spaces-divisors-]{spaces-divisors}
\externaldocument[spaces-over-fields-]{spaces-over-fields}
\externaldocument[spaces-topologies-]{spaces-topologies}
\externaldocument[spaces-descent-]{spaces-descent}
\externaldocument[spaces-perfect-]{spaces-perfect}
\externaldocument[spaces-more-morphisms-]{spaces-more-morphisms}
\externaldocument[spaces-flat-]{spaces-flat}
\externaldocument[spaces-groupoids-]{spaces-groupoids}
\externaldocument[spaces-more-groupoids-]{spaces-more-groupoids}
\externaldocument[bootstrap-]{bootstrap}
\externaldocument[spaces-pushouts-]{spaces-pushouts}
\externaldocument[spaces-chow-]{spaces-chow}
\externaldocument[groupoids-quotients-]{groupoids-quotients}
\externaldocument[spaces-more-cohomology-]{spaces-more-cohomology}
\externaldocument[spaces-simplicial-]{spaces-simplicial}
\externaldocument[spaces-duality-]{spaces-duality}
\externaldocument[formal-spaces-]{formal-spaces}
\externaldocument[restricted-]{restricted}
\externaldocument[spaces-resolve-]{spaces-resolve}
\externaldocument[formal-defos-]{formal-defos}
\externaldocument[defos-]{defos}
\externaldocument[cotangent-]{cotangent}
\externaldocument[examples-defos-]{examples-defos}
\externaldocument[algebraic-]{algebraic}
\externaldocument[examples-stacks-]{examples-stacks}
\externaldocument[stacks-sheaves-]{stacks-sheaves}
\externaldocument[criteria-]{criteria}
\externaldocument[artin-]{artin}
\externaldocument[quot-]{quot}
\externaldocument[stacks-properties-]{stacks-properties}
\externaldocument[stacks-morphisms-]{stacks-morphisms}
\externaldocument[stacks-limits-]{stacks-limits}
\externaldocument[stacks-cohomology-]{stacks-cohomology}
\externaldocument[stacks-perfect-]{stacks-perfect}
\externaldocument[stacks-introduction-]{stacks-introduction}
\externaldocument[stacks-more-morphisms-]{stacks-more-morphisms}
\externaldocument[stacks-geometry-]{stacks-geometry}
\externaldocument[moduli-]{moduli}
\externaldocument[moduli-curves-]{moduli-curves}
\externaldocument[examples-]{examples}
\externaldocument[exercises-]{exercises}
\externaldocument[guide-]{guide}
\externaldocument[desirables-]{desirables}
\externaldocument[coding-]{coding}
\externaldocument[obsolete-]{obsolete}
\externaldocument[fdl-]{fdl}
\externaldocument[index-]{index}

% Theorem environments.
%
\theoremstyle{plain}
\newtheorem{theorem}[subsection]{Theorem}
\newtheorem{proposition}[subsection]{Proposition}
\newtheorem{lemma}[subsection]{Lemma}

\theoremstyle{definition}
\newtheorem{definition}[subsection]{Definition}
\newtheorem{example}[subsection]{Example}
\newtheorem{exercise}[subsection]{Exercise}
\newtheorem{situation}[subsection]{Situation}

\theoremstyle{remark}
\newtheorem{remark}[subsection]{Remark}
\newtheorem{remarks}[subsection]{Remarks}

\numberwithin{equation}{subsection}

% Macros
%
\def\lim{\mathop{\mathrm{lim}}\nolimits}
\def\colim{\mathop{\mathrm{colim}}\nolimits}
\def\Spec{\mathop{\mathrm{Spec}}}
\def\Hom{\mathop{\mathrm{Hom}}\nolimits}
\def\Ext{\mathop{\mathrm{Ext}}\nolimits}
\def\SheafHom{\mathop{\mathcal{H}\!\mathit{om}}\nolimits}
\def\SheafExt{\mathop{\mathcal{E}\!\mathit{xt}}\nolimits}
\def\Sch{\mathit{Sch}}
\def\Mor{\mathop{\mathrm{Mor}}\nolimits}
\def\Ob{\mathop{\mathrm{Ob}}\nolimits}
\def\Sh{\mathop{\mathit{Sh}}\nolimits}
\def\NL{\mathop{N\!L}\nolimits}
\def\CH{\mathop{\mathrm{CH}}\nolimits}
\def\proetale{{pro\text{-}\acute{e}tale}}
\def\etale{{\acute{e}tale}}
\def\QCoh{\mathit{QCoh}}
\def\Ker{\mathop{\mathrm{Ker}}}
\def\Im{\mathop{\mathrm{Im}}}
\def\Coker{\mathop{\mathrm{Coker}}}
\def\Coim{\mathop{\mathrm{Coim}}}

% Boxtimes
%
\DeclareMathSymbol{\boxtimes}{\mathbin}{AMSa}{"02}

%
% Macros for moduli stacks/spaces
%
\def\QCohstack{\mathcal{QC}\!\mathit{oh}}
\def\Cohstack{\mathcal{C}\!\mathit{oh}}
\def\Spacesstack{\mathcal{S}\!\mathit{paces}}
\def\Quotfunctor{\mathrm{Quot}}
\def\Hilbfunctor{\mathrm{Hilb}}
\def\Curvesstack{\mathcal{C}\!\mathit{urves}}
\def\Polarizedstack{\mathcal{P}\!\mathit{olarized}}
\def\Complexesstack{\mathcal{C}\!\mathit{omplexes}}
% \Pic is the operator that assigns to X its picard group, usage \Pic(X)
% \Picardstack_{X/B} denotes the Picard stack of X over B
% \Picardfunctor_{X/B} denotes the Picard functor of X over B
\def\Pic{\mathop{\mathrm{Pic}}\nolimits}
\def\Picardstack{\mathcal{P}\!\mathit{ic}}
\def\Picardfunctor{\mathrm{Pic}}
\def\Deformationcategory{\mathcal{D}\!\mathit{ef}}
\usepackage{cleveref}
\crefname{originaldefinition}{original definition}{original definitions}
\Crefname{originaldefinition}{Original definition}{Original definitions}
\makeatletter
\newlabel{situation-complexes}{{0DLE}{source}{}{sourcepointer-1}{}}
\newlabel{situation-complexes@cref}{{[originaldefinition][0][]0DLE}{[1][0][]source}{}{}{}}
\makeatother
\crefname{originallemma}{original lemma}{original lemmas}
\Crefname{originallemma}{Original lemma}{Original lemmas}
\makeatletter
\newlabel{lemma-complexes-diagonal}{{0DLG}{source}{}{sourcepointer-2}{}}
\newlabel{lemma-complexes-diagonal@cref}{{[originallemma][0][]0DLG}{[1][0][]source}{}{}{}}
\makeatother
\crefname{originallemma}{original lemma}{original lemmas}
\Crefname{originallemma}{Original lemma}{Original lemmas}
\makeatletter
\newlabel{algebraic-lemma-representable-diagonal}{{045G}{source}{}{sourcepointer-3}{}}
\newlabel{algebraic-lemma-representable-diagonal@cref}{{[originallemma][0][]045G}{[1][0][]source}{}{}{}}
\makeatother
\crefname{originallemma}{original lemma}{original lemmas}
\Crefname{originallemma}{Original lemma}{Original lemmas}
\makeatletter
\newlabel{spaces-descent-lemma-fpqc-universal-effective-epimorphisms}{{04P2}{source}{}{sourcepointer-4}{}}
\newlabel{spaces-descent-lemma-fpqc-universal-effective-epimorphisms@cref}{{[originallemma][0][]04P2}{[1][0][]source}{}{}{}}
\makeatother
\crefname{originallemma}{original lemma}{original lemmas}
\Crefname{originallemma}{Original lemma}{Original lemmas}
\makeatletter
\newlabel{lemma-complexes}{{0DLD}{source}{}{sourcepointer-5}{}}
\newlabel{lemma-complexes@cref}{{[originallemma][0][]0DLD}{[1][0][]source}{}{}{}}
\makeatother
\crefname{originallemma}{original lemma}{original lemmas}
\Crefname{originallemma}{Original lemma}{Original lemmas}
\makeatletter
\newlabel{spaces-simplicial-lemma-fppf-glue-neg-ext-zero}{{0DLA}{source}{}{sourcepointer-6}{}}
\newlabel{spaces-simplicial-lemma-fppf-glue-neg-ext-zero@cref}{{[originallemma][0][]0DLA}{[1][0][]source}{}{}{}}
\makeatother
\crefname{originallemma}{original lemma}{original lemmas}
\Crefname{originallemma}{Original lemma}{Original lemmas}
\makeatletter
\newlabel{spaces-more-morphisms-lemma-pseudo-coherent-descends-fpqc}{{0DL1}{source}{}{sourcepointer-7}{}}
\newlabel{spaces-more-morphisms-lemma-pseudo-coherent-descends-fpqc@cref}{{[originallemma][0][]0DL1}{[1][0][]source}{}{}{}}
\makeatother
\crefname{originallemma}{original lemma}{original lemmas}
\Crefname{originallemma}{Original lemma}{Original lemmas}
\makeatletter
\newlabel{spaces-more-morphisms-lemma-tor-amplitude-descends-fppf}{{0DL2}{source}{}{sourcepointer-8}{}}
\newlabel{spaces-more-morphisms-lemma-tor-amplitude-descends-fppf@cref}{{[originallemma][0][]0DL2}{[1][0][]source}{}{}{}}
\makeatother
\crefname{originallemma}{original lemma}{original lemmas}
\Crefname{originallemma}{Original lemma}{Original lemmas}
\makeatletter
\newlabel{lemma-complexes-open-neg-exts-vanishing}{{0DLC}{source}{}{sourcepointer-9}{}}
\newlabel{lemma-complexes-open-neg-exts-vanishing@cref}{{[originallemma][0][]0DLC}{[1][0][]source}{}{}{}}
\makeatother
\title{The stack of relatively perfect complexes is a stack in groupoids}
\author{The Stacks Project Authors}
\date{stacks-project:quot.tex:local-snapshot-20261012; GNU FDL 1.2 or later}
\begin{document}
\maketitle
\setcounter{subsection}{1}
\noindent\textit{Editorial note.} All delivered mathematical text is the author's own, copied from the pinned Stacks Project chapter file quot.tex; no statement or proof step was rewritten, reordered or removed. The counted span is the author environment and its own proof environment. References inside the span are handed to the assembler as navigation stubs that carry the canonical tag of the cited tag only; whether each of them is load-bearing is recorded on the stub itself.\par\medskip
\noindent\textit{Allowed standard prerequisites (editorial).}\begin{itemize}
\item Situation situation-complexes (canonical tag 0DLE). cited at statement line 5493. Standing situation: X -> B and the relatively perfect K.
\item Lemma lemma-complexes-diagonal (canonical tag 0DLG). cited at proof line 5502. The diagonal of the complexes stack is representable.
\item Lemma algebraic-lemma-representable-diagonal (canonical tag 045G). cited at proof lines 5502-5503. Representable diagonal makes Isom a sheaf.
\item Lemma spaces-descent-lemma-fpqc-universal-effective-epimorphisms (canonical tag 04P2). cited at proof lines 5517-5518. The morphisms g\_i glue uniquely to g.
\item Lemma lemma-complexes (canonical tag 0DLD). cited at proof lines 5539-5540 and 5559. Negative Ext vanishing enables descent; condition (2) characterized.
\item Lemma spaces-simplicial-lemma-fppf-glue-neg-ext-zero (canonical tag 0DLA). cited at proof lines 5541 and 5548 (footnote). Descent of complexes along the Cech fppf hypercovering.
\item Lemma spaces-more-morphisms-lemma-pseudo-coherent-descends-fpqc (canonical tag 0DL1). cited at proof lines 5555-5556. Pseudo-coherence of the descended complex.
\item Lemma spaces-more-morphisms-lemma-tor-amplitude-descends-fppf (canonical tag 0DL2). cited at proof lines 5556-5557. Tor amplitude descends.
\item Lemma lemma-complexes-open-neg-exts-vanishing (canonical tag 0DLC). cited at proof lines 5560-5561. Open W verifies condition (2) for the descended object.
\end{itemize}\medskip
\ifdefined\hypertarget\hypertarget{sourcepointer-1}{}\fi
\noindent\textit{Source reference.} Original definition 0DLE: sources/stacks-chow/tags.txt, source lines 17044--17044. Cited at statement line 5493: the lemma is phrased in this situation; load-bearing.\par\medskip
\ifdefined\hypertarget\hypertarget{sourcepointer-2}{}\fi
\noindent\textit{Source reference.} Original lemma 0DLG: sources/stacks-chow/tags.txt, source lines 17046--17046. Cited at proof lines 5502: Isom presheaves are sheaves via the diagonal; load-bearing.\par\medskip
\ifdefined\hypertarget\hypertarget{sourcepointer-3}{}\fi
\noindent\textit{Source reference.} Original lemma 045G: sources/stacks-chow/tags.txt, source lines 5435--5435. Cited at proof lines 5502-5503: the transfer lemma from diagonal to Isom sheaves; load-bearing.\par\medskip
\ifdefined\hypertarget\hypertarget{sourcepointer-4}{}\fi
\noindent\textit{Source reference.} Original lemma 04P2: sources/stacks-chow/tags.txt, source lines 6100--6100. Cited at proof lines 5517-5518: g\_i glue uniquely to g : T -> B; load-bearing.\par\medskip
\ifdefined\hypertarget\hypertarget{sourcepointer-5}{}\fi
\noindent\textit{Source reference.} Original lemma 0DLD: sources/stacks-chow/tags.txt, source lines 17043--17043. Cited at proof lines 5539-5540 (vanishing enabling descent) and 5559 (condition (2) for the descended E); load-bearing twice.\par\medskip
\ifdefined\hypertarget\hypertarget{sourcepointer-6}{}\fi
\noindent\textit{Source reference.} Original lemma 0DLA: sources/stacks-chow/tags.txt, source lines 17040--17040. Cited at proof lines 5541 and 5548 (footnote Cech hypercovering): descent for the complexes; load-bearing twice.\par\medskip
\ifdefined\hypertarget\hypertarget{sourcepointer-7}{}\fi
\noindent\textit{Source reference.} Original lemma 0DL1: sources/stacks-chow/tags.txt, source lines 17031--17031. Cited at proof lines 5555-5556: E is pseudo-coherent; load-bearing.\par\medskip
\ifdefined\hypertarget\hypertarget{sourcepointer-8}{}\fi
\noindent\textit{Source reference.} Original lemma 0DL2: sources/stacks-chow/tags.txt, source lines 17032--17032. Cited at proof lines 5556-5557: locally finite tor dimension descends, so E is T-perfect; load-bearing.\par\medskip
\ifdefined\hypertarget\hypertarget{sourcepointer-9}{}\fi
\noindent\textit{Source reference.} Original lemma 0DLC: sources/stacks-chow/tags.txt, source lines 17042--17042. Cited at proof lines 5560-5561: the description of the open W verifies condition (2) for E; load-bearing.\par\medskip
\begin{lemma}
\label{lemma-complexes-stack}
In Situation \ref{situation-complexes} the functor
$p : \Complexesstack_{X/B} \longrightarrow (\Sch/S)_{fppf}$
is a stack in groupoids.
\end{lemma}

\begin{proof}
To prove that $\Complexesstack_{X/B}$ is a stack in groupoids,
we have to show that the presheaves $\mathit{Isom}$ are sheaves
and that descent data are effective. The statement on
$\mathit{Isom}$ follows from Lemma \ref{lemma-complexes-diagonal}, see
Algebraic Stacks, Lemma \ref{algebraic-lemma-representable-diagonal}.
Let us prove the statement on descent data.

\medskip\noindent
Suppose that $\{a_i : T_i \to T\}$ is an fppf covering of schemes over $S$.
Let $(\xi_i, \varphi_{ij})$ be a descent datum for $\{T_i \to T\}$
with values in $\Complexesstack_{X/B}$.
For each $i$ we can write $\xi_i = (T_i, g_i, E_i)$.
Denote $\text{pr}_0 : T_i \times_T T_j \to T_i$ and
$\text{pr}_1 : T_i \times_T T_j \to T_j$ the projections.
The condition that $\xi_i|_{T_i \times_T T_j} \cong \xi_j|_{T_i \times_T T_j}$
implies in particular that $g_i \circ \text{pr}_0 = g_j \circ \text{pr}_1$.
Thus there exists a unique morphism $g : T \to B$ such that
$g_i = g \circ a_i$, see
Descent on Spaces, Lemma
\ref{spaces-descent-lemma-fpqc-universal-effective-epimorphisms}.
Denote $X_T = T \times_{g, B} X$. Set
$X_i = X_{T_i} = T_i \times_{g_i, B} X = T_i \times_{a_i, T} X_T$
and
$$
X_{ij} = X_{T_i} \times_{X_T} X_{T_j} = X_i \times_{X_T} X_j
$$
with projections $\text{pr}_i$ and $\text{pr}_j$ to $X_i$ and $X_j$.
Observe that the pullback of $(T_i, g_i, E_i)$
by $\text{pr}_0 : T_i \times_T T_j \to T_i$ is given by
$(T_i \times_T T_j, g_i \circ \text{pr}_0, L\text{pr}_i^*E_i)$.
Hence a descent datum for $\{T_i \to T\}$ in $\Complexesstack_{X/B}$
is given by the objects $(T_i, g \circ a_i, E_i)$
and for each pair $i, j$ an isomorphism in
$D\mathcal{O}_{X_{ij}})$
$$
\varphi_{ij} :
L\text{pr}_i^*E_i \longrightarrow L\text{pr}_j^*E_j
$$
satisfying the cocycle condition over the pullback of $X$ to
$T_i \times_T T_j \times_T T_k$.
Using the vanishing of negative Exts provided by (b) of
Lemma \ref{lemma-complexes}, we may apply
Simplicial Spaces, Lemma \ref{spaces-simplicial-lemma-fppf-glue-neg-ext-zero}
to obtain descent\footnote{To check this, first consider the
{\v C}ech fppf hypercovering $a : T_\bullet \to T$ with
$T_n = \coprod_{i_0 \ldots i_n} T_{i_0} \times_T \ldots \times_T T_{i_n}$
and then set $U_\bullet = T_\bullet \times_{a, T} X_T$. Then
$U_\bullet \to X_T$ is an fppf hypercovering to which we may
apply Simplicial Spaces, Lemma
\ref{spaces-simplicial-lemma-fppf-glue-neg-ext-zero}.}
for these complexes. In other words, we find there exists an object
$E$ in $D_\QCoh(\mathcal{O}_{X_T})$ restricting to $E_i$ on $X_{T_i}$
compatible with $\varphi_{ij}$. Recall that being
$T$-perfect signifies being pseudo-coherent and having
locally finite tor dimension over $f^{-1}\mathcal{O}_T$.
Thus $E$ is $T$-perfect by an application of
More on Morphisms of Spaces, Lemmas
\ref{spaces-more-morphisms-lemma-pseudo-coherent-descends-fpqc} and
\ref{spaces-more-morphisms-lemma-tor-amplitude-descends-fppf}.
Finally, we have to check condition (2) from
Lemma \ref{lemma-complexes} for $E$.
This immediately follows from the description of the open $W$
in Lemma \ref{lemma-complexes-open-neg-exts-vanishing}
and the fact that (2) holds for $E_i$ on $X_{T_i}/T_i$.
\end{proof}

\end{document}
